% GENERATED FILE. Edit canonical lessons, then run npm run generate:books.
% canonical-source-hash: fnv1a64:34cb2c101564cb75
% canonical-lessons: MR-C289-saksidar, MR-C289-gunha, MR-C289-daroda, MR-C289-pakitmar, MR-C289-phasavne

\chapter{Crimes and Witnesses}
\label{ch:mr-crimes-and-witnesses}
\hlchaptermodality{289}

\begin{chapteropening}
\textbf{By the end of this chapter:} \emph{I can talk about a crime, a robbery and the witness who saw it.}

Complete the last lesson of chapter 289: I can talk about a crime, a robbery and the witness who saw it.
\end{chapteropening}

\section[\texorpdfstring{\mr{साक्षीदार} (sākṣīdār)}{साक्षीदार (sākṣīdār)}]{\mr{साक्षीदार} (sākṣīdār) — a witness}

\label{lesson:MR-C289-saksidar}



\begin{quote}
Before the new one: say the Marathi for to sink, then the Marathi for minor.
\end{quote}

\subsection*{You'll want to know: \mr{साक्षीदार}}
\textbf{\mr{साक्षीदार}} — \emph{sākṣīdār} — \textquotedblleft{}a witness\textquotedblright{}.

\begin{quote}
\mr{त्या} \mr{साक्षीदाराने} \mr{सगळं} \mr{पाहिलं} \mr{असतं} \mr{तर} \mr{त्याने} \mr{पोलिसांना} \mr{सांगितलं} \mr{असतं}.
\end{quote}

\begin{quote}
\emph{tyā sākṣīdārāne sagḷaṁ pāhilaṁ astaṁ tar tyāne polisānnā sāṅgitlaṁ astaṁ.}
\end{quote}

\begin{quote}
If that witness had seen everything, he would have told the police.
\end{quote}

Masculine.

\mr{साक्षीदाराने}, by the witness; a past form with \mr{असतं} says what would have happened.

\begin{cousinweb}
\textbf{\mr{साक्षीदार}} (\emph{sākṣīdār}) is built from Sanskrit \emph{sākṣin}, 'an eyewitness' (Monier-Williams), and the Persian suffix \emph{-dār}, 'one who holds': the one who holds the testimony.
\end{cousinweb}

\begin{grammarlens}[title={What you've built}]
One more word for telling what happened.
\end{grammarlens}

\subsection*{Your turn}
\begin{itemize}
\raggedright
  \item \emph{Say it:} \emph{sākṣīdār}
  \item \emph{Say it:} \emph{sākṣīdār}, once more
  \item \emph{Say it:} \emph{tyā sākṣīdārāne sagḷaṁ pāhilaṁ astaṁ tar tyāne polisānnā sāṅgitlaṁ astaṁ.} — If that witness had seen everything, he would have told the police.
  \item \emph{Recall:} say the Marathi for to sink, then the Marathi for minor, then say \emph{sākṣīdār} again
\end{itemize}

\begin{tcolorbox}[breakable,colback=teal!4,colframe=teal!35!black,title={Before you move on}]
What does \mr{साक्षीदार} mean? (\textquotedblleft{}A witness\textquotedblright{}.) Say it once more.
\end{tcolorbox}
\section[\texorpdfstring{\mr{गुन्हा} (gunhā)}{गुन्हा (gunhā)}]{\mr{गुन्हा} (gunhā) — a crime; an offence}

\label{lesson:MR-C289-gunha}



\begin{quote}
Before the new one: say the Marathi for minor, then the Marathi for a witness.
\end{quote}

\subsection*{You'll want to know: \mr{गुन्हा}}
\textbf{\mr{गुन्हा}} — \emph{gunhā} — \textquotedblleft{}a crime; an offence\textquotedblright{}.

\begin{quote}
\mr{खोटी} \mr{सही} \mr{करणं} \mr{हा} \mr{मोठा} \mr{गुन्हा} \mr{आहे}.
\end{quote}

\begin{quote}
\emph{khoṭī sahī karṇaṁ hā moṭhā gunhā āhe.}
\end{quote}

\begin{quote}
Forging a signature is a serious crime.
\end{quote}

Masculine.

The oblique is \mr{गुन्ह्या}-: \mr{गुन्ह्यात}, in the crime.

\begin{cousinweb}
\textbf{\mr{गुन्हा}} (\emph{gunhā}) is a loan from Persian \emph{gunāh}, 'a sin, a fault'; in Marathi it became the everyday and the legal word for an offence.
\end{cousinweb}

\begin{grammarlens}[title={What you've built}]
One more word for telling what happened.
\end{grammarlens}

\subsection*{Your turn}
\begin{itemize}
\raggedright
  \item \emph{Say it:} \emph{gunhā}
  \item \emph{Say it:} \emph{gunhā}, once more
  \item \emph{Say it:} \emph{khoṭī sahī karṇaṁ hā moṭhā gunhā āhe.} — Forging a signature is a serious crime.
  \item \emph{Recall:} say the Marathi for minor, then the Marathi for a witness, then say \emph{gunhā} again
\end{itemize}

\begin{tcolorbox}[breakable,colback=teal!4,colframe=teal!35!black,title={Before you move on}]
What does \mr{गुन्हा} mean? (\textquotedblleft{}A crime; an offence\textquotedblright{}.) Say it once more.
\end{tcolorbox}
\section[\texorpdfstring{\mr{दरोडा} (daroḍā)}{दरोडा (daroḍā)}]{\mr{दरोडा} (daroḍā) — a robbery; a gang raid}

\label{lesson:MR-C289-daroda}



\begin{quote}
Before the new one: say the Marathi for a witness, then the Marathi for a crime.
\end{quote}

\subsection*{You'll want to know: \mr{दरोडा}}
\textbf{\mr{दरोडा}} — \emph{daroḍā} — \textquotedblleft{}a robbery; a gang raid\textquotedblright{}.

\begin{quote}
\mr{रात्री} \mr{चार} \mr{माणसांनी} \mr{दुकानावर} \mr{दरोडा} \mr{टाकला}.
\end{quote}

\begin{quote}
\emph{rātrī cār māṇsānnī dukānāvar daroḍā ṭāklā.}
\end{quote}

\begin{quote}
At night four men raided the shop.
\end{quote}

Masculine.

\mr{दरोडा} \mr{पडणे}, a robbery takes place; \mr{दरोडा} \mr{टाकणे}, to carry one out.

\begin{cousinweb}
\textbf{\mr{दरोडा}} (\emph{daroḍā}) is inherited: CDIAL lists Marathi \emph{daroḍā}, 'an attack by a band of robbers', under \emph{*daḍavaḍa}, 'an onrush' (CDIAL 6127).
\end{cousinweb}

\begin{grammarlens}[title={What you've built}]
One more word for telling what happened.
\end{grammarlens}

\subsection*{Your turn}
\begin{itemize}
\raggedright
  \item \emph{Say it:} \emph{daroḍā}
  \item \emph{Say it:} \emph{daroḍā}, once more
  \item \emph{Say it:} \emph{rātrī cār māṇsānnī dukānāvar daroḍā ṭāklā.} — At night four men raided the shop.
  \item \emph{Recall:} say the Marathi for a witness, then the Marathi for a crime, then say \emph{daroḍā} again
\end{itemize}

\begin{tcolorbox}[breakable,colback=teal!4,colframe=teal!35!black,title={Before you move on}]
What does \mr{दरोडा} mean? (\textquotedblleft{}A robbery; a gang raid\textquotedblright{}.) Say it once more.
\end{tcolorbox}
\section[\texorpdfstring{\mr{पाकीटमार} (pākīṭmār)}{पाकीटमार (pākīṭmār)}]{\mr{पाकीटमार} (pākīṭmār) — a pickpocket}

\label{lesson:MR-C289-pakitmar}



\begin{quote}
Before the new one: say the Marathi for a crime, then the Marathi for a robbery.
\end{quote}

\subsection*{You'll want to know: \mr{पाकीटमार}}
\textbf{\mr{पाकीटमार}} — \emph{pākīṭmār} — \textquotedblleft{}a pickpocket\textquotedblright{}.

\begin{quote}
\mr{गर्दीत} \mr{एका} \mr{पाकीटमाराने} \mr{माझं} \mr{पाकीट} \mr{मारलं}.
\end{quote}

\begin{quote}
\emph{gardīt ekā pākīṭmārāne mājhaṁ pākīṭ mārlaṁ.}
\end{quote}

\begin{quote}
In the crowd a pickpocket stole my wallet.
\end{quote}

Masculine.

Built from \mr{पाकीट}, a wallet, and \mr{मार}-, from \mr{मारणे}: \mr{पाकीट} \mr{मारणे}, to pick a pocket.

\begin{cousinweb}
\textbf{\mr{पाकीटमार}} (\emph{pākīṭmār}) is a Marathi compound built from \emph{pākīṭ}, 'a wallet', a loan from English \emph{packet}, and \emph{mār-}, the stem of \emph{mārṇe}, 'to strike' (CDIAL 10066).
\end{cousinweb}

\begin{grammarlens}[title={What you've built}]
One more word for telling what happened.
\end{grammarlens}

\subsection*{Your turn}
\begin{itemize}
\raggedright
  \item \emph{Say it:} \emph{pākīṭmār}
  \item \emph{Say it:} \emph{pākīṭmār}, once more
  \item \emph{Say it:} \emph{gardīt ekā pākīṭmārāne mājhaṁ pākīṭ mārlaṁ.} — In the crowd a pickpocket stole my wallet.
  \item \emph{Recall:} say the Marathi for a crime, then the Marathi for a robbery, then say \emph{pākīṭmār} again
\end{itemize}

\begin{tcolorbox}[breakable,colback=teal!4,colframe=teal!35!black,title={Before you move on}]
What does \mr{पाकीटमार} mean? (\textquotedblleft{}A pickpocket\textquotedblright{}.) Say it once more.
\end{tcolorbox}
\section[\texorpdfstring{\mr{फसवणे} (phasavṇe)}{फसवणे (phasavṇe)}]{\mr{फसवणे} (phasavṇe) — to cheat; to deceive}

\label{lesson:MR-C289-phasavne}



\begin{quote}
Before the new one: say the Marathi for a robbery, then the Marathi for a pickpocket.
\end{quote}

\subsection*{You'll want to know: \mr{फसवणे}}
\textbf{\mr{फसवणे}} — \emph{phasavṇe} — \textquotedblleft{}to cheat; to deceive\textquotedblright{}.

\begin{quote}
\mr{त्या} \mr{दुकानदाराने} \mr{आम्हाला} \mr{फसवलं}.
\end{quote}

\begin{quote}
\emph{tyā dukāndārāne āmhālā phasavlaṁ.}
\end{quote}

\begin{quote}
That shopkeeper cheated us.
\end{quote}

Past, with -\mr{ने} (he …): \mr{त्याने} \mr{फसवलं}, \emph{tyāne phasavlaṁ}.

The one cheated takes -\mr{ला}: \mr{लोकांना} \mr{फसवणे}, to cheat people.

\begin{cousinweb}
\textbf{\mr{फसवणे}} (\emph{phasavṇe}) is inherited: CDIAL lists Marathi \emph{phasavṇe}, 'to deceive', under Sanskrit \emph{spāśayati}, 'causes to be bound or stopped' (CDIAL 13814). To cheat is to snare.
\end{cousinweb}

\begin{grammarlens}[title={What you've built}]
One more word for telling what happened.
\end{grammarlens}

\subsection*{Your turn}
\begin{itemize}
\raggedright
  \item \emph{Say it:} \emph{phasavṇe}
  \item \emph{Say it:} \emph{phasavṇe}, once more
  \item \emph{Say it:} \emph{tyā dukāndārāne āmhālā phasavlaṁ.} — That shopkeeper cheated us.
  \item \emph{Recall:} say the Marathi for a robbery, then the Marathi for a pickpocket, then say \emph{phasavṇe} again
\end{itemize}

\begin{tcolorbox}[breakable,colback=teal!4,colframe=teal!35!black,title={Before you move on}]
What does \mr{फसवणे} mean? (\textquotedblleft{}To cheat; to deceive\textquotedblright{}.) Say it once more.
\end{tcolorbox}
